%=======================================================================
%  IE 5311-001 (D01) Exploration Problem EP5
%  One page, single spacing, 10 pt, 1-inch margins, as the sheet requires.
%  Content past page 1 is not graded, so the page limit is hard.
%  Sections map onto the two parts asked: the main idea of PDHG, and
%  how cuPDLP applies it to large-scale linear programs.
%
%  IF IT SPILLS ONTO A SECOND PAGE, delete in this order until it fits.
%    1. the last sentence of section 5 (from "On small and medium")
%    2. the "Presolve and scaling" row of the table in section 3
%    3. the last sentence of section 1
%  Do not cut section 2, section 4, or the restart row of the table:
%  they are what the question asks for.
%=======================================================================
\documentclass[10pt]{article}

\usepackage[margin=1in]{geometry}
\usepackage{amsmath}
\usepackage{amssymb}
\usepackage{booktabs}
\usepackage{array}

\newcommand{\proj}{\mathrm{proj}}

\pagestyle{empty}
\setlength{\parskip}{0.35em}
\setlength{\parindent}{0pt}

\begin{document}

{\large\bfseries EP5.\quad PDHG and cuPDLP for large-scale linear programs}\\[1pt]
{\small Lidivine Kengne \quad\textbar\quad IE 5311-001 (D01), Homework 5}
\vspace{0.3em}

\textbf{1. Why a first-order method for LP at all.}
Both mature methods solve linear systems. Simplex keeps a
factorization of the basis $A_B$ and updates it; the interior point
method forms and factors $A \Theta A^{\!\top}$ every iteration. That
is where the memory goes: even for sparse $A$ the factors fill in, and
the fill is not under the modeller's control. At hundreds of millions
of nonzeros the factors stop fitting, and the solver fails for reasons
unrelated to how many iterations it would have needed.

\textbf{2. The main idea of PDHG.}
Take $\min c^{\!\top}x$ subject to $Ax \ge b$, $x \ge 0$. Its
Lagrangian, with $y \ge 0$ on the rows, is
$L(x,y) = c^{\!\top}x + y^{\!\top}(b - Ax)$, and strong duality makes
the LP equivalent to the saddle point problem
$\min_{x \ge 0} \max_{y \ge 0} L(x,y)$. This is the Lagrangian derived
in Problem~2 Part~5; the difference is that there $x$ was minimised out
to leave a dual, whereas here \emph{neither} side is eliminated.
Primal--dual hybrid gradient takes one projected gradient step in each
variable, extrapolating the primal:
\[
  x^{k+1} = \proj_{x \ge 0}\!\bigl(x^{k} - \tau (c - A^{\!\top} y^{k})\bigr),
  \quad
  y^{k+1} = \proj_{y \ge 0}\!\bigl(y^{k} + \sigma \bigl(b - A(2x^{k+1} - x^{k})\bigr)\bigr).
\]
Projection onto an orthant or a box is a componentwise $\max$, so an
iteration costs exactly one $Ax$ and one $A^{\!\top}y$:
$O(\mathrm{nnz}(A))$ work and $O(\mathrm{nnz}(A))$ memory, no fill-in
ever, no linear system anywhere. Convergence needs
$\tau\sigma\|A\|_2^2 < 1$ and the averaged iterates close the
primal--dual gap at rate $O(1/k)$. The term $2x^{k+1}-x^{k}$ is not
cosmetic: without it this is Arrow--Hurwicz, which need not converge.

\textbf{3. From PDHG to PDLP.}
Textbook PDHG on a real LP stalls at one or two digits, so the
algorithm above is necessary but nowhere near sufficient. PDLP is the
engineering that closes the gap:

\begin{center}\small
\begin{tabular}{@{}l p{10.2cm}@{}}
\toprule
Restarting & Restart periodically from the running average. On an LP
the normalised duality gap is \emph{sharp}, which turns the $O(1/k)$
rate into a linear one. The single largest gain.\\
Primal weight & Set $\tau = \eta/\omega$, $\sigma = \eta\omega$ and
adapt $\omega$, so primal and dual residuals fall together.\\
Adaptive step & Replace the global, pessimistic bound
$\eta < 1/\|A\|_2$ by a step accepted or rejected on local curvature.\\
Presolve and scaling & Ruiz equilibration and diagonal
preconditioning; conditioning of $A$ sets the constant in every bound
above.\\
\bottomrule
\end{tabular}
\end{center}

\textbf{4. How cuPDLP applies this at scale.}
cuPDLP is PDLP moved onto a GPU, and the reason it works is that the
iteration was already the right shape. Each step is two sparse
matrix--vector products plus componentwise arithmetic and a few
reductions for the restart and step-size tests --- all data-parallel,
all bandwidth-bound, which is precisely what GPU memory systems
deliver. Nothing in the method is sequential: there is no pivot order
to respect, no factorization to build, no triangular solve. That is
the contrast with the alternatives, since a simplex pivot sequence is
inherently serial and a sparse Cholesky parallelises only within its
elimination tree while filling in badly. Memory tells the same story:
the method holds $A$ in sparse form plus a handful of vectors of
length $n$ and $m$, so instances whose \emph{factors} would exhaust
host memory have a \emph{footprint} that fits on a single device.
Large and sparse is therefore the regime where the method is
strongest, and it is why first-order LP moved from curiosity to
practice and has since appeared in commercial solvers.

\textbf{5. What it costs.}
Accuracy: relative gaps of $10^{-4}$ to $10^{-8}$ are routine and
$10^{-9}$ is hard, where simplex returns an exact vertex. No basis:
hence no warm start for branch and bound and no clean sensitivity
analysis, the duals coming out approximate --- which matters for
exactly the shadow-price reading done in Homework~1. And
predictability: performance tracks conditioning in a way that is hard
to forecast from the model. On small and medium LPs, where a
factorization is cheap, simplex and interior point remain correct; the
case for PDHG begins where the factors stop fitting.

\end{document}
